\chapter{Tiempo de Revisión y Análisis de Errores (OE3)}
\label{cap:utilidad}

% ===================================================================
\section{Introducción}
% ===================================================================

Este capítulo corresponde al tercer objetivo específico:

\begin{quote}
    \textbf{OE3.} Cuantificar la reducción del tiempo de revisión que aporta el
    sistema y caracterizar sus errores mediante revisión experta de una muestra de
    predicciones.
\end{quote}

Según el cronograma (Tabla~\ref{tab:cronograma}), sus tres resultados se
entregan en E3; en esta entrega el capítulo fija su estructura. Los resultados
comprometidos en la Tabla~\ref{tab:iov-oe3} ocuparán una sección cada uno: la
metodología de validación comparativa (RE3.1, sección~\ref{sec:re31}), la
medición comparativa del tiempo de revisión (RE3.2, sección~\ref{sec:re32}) y el
análisis cualitativo de errores y hallazgos (RE3.3, sección~\ref{sec:re33}).
Cada sección seguirá la estructura de los capítulos~\ref{cap:curacion}
y~\ref{cap:deteccion}, y la sección~\ref{sec:discusion-oe3} discutirá los tres
en conjunto.

El material de entrada es el modelo seleccionado en el
Capítulo~\ref{cap:deteccion} y la partición de prueba construida en el
Capítulo~\ref{cap:curacion}: OE3 no entrena ni vuelve a ajustar nada, mide lo que
cuesta usar el sistema.

% ===================================================================
\section{Metodología de validación comparativa (RE3.1)}
\label{sec:re31}
% ===================================================================
% TODO Redactar. El documento de metodología se fecha ANTES de ejecutar la
% evaluación; fijarlo de antemano impide elegir a posteriori el umbral o la
% métrica que mejor favorezcan al resultado ya obtenido.
%
% El IOV comprometido (Tabla~\ref{tab:iov-oe3}) obliga a especificar, en este
% orden:
%   (1) el subconjunto de prueba, particionado por archivo de grabación y sin
%       solapamiento con entrenamiento ni validación;
%   (2) el punto de operación: confianza 0,5 y NMS IoU 0,3;
%   (3) el umbral IoU 0,5 para contar un verdadero positivo;
%   (4) las métricas que se reportan;
%   (5) el emparejamiento entre predicción y anotación;
%   (6) el procedimiento para revisar discrepancias y eventos potencialmente
%       omitidos por la anotación de referencia.
%
% ATENCIÓN, CONFLICTO A RESOLVER ANTES DE REDACTAR.
% El protocolo implementado (src/evaluation/protocol.py, descrito en
% capitulos/05_oe2_detector/protocolo_comparacion.md) NO usa este punto de
% operación: empareja a IoU 0,3 y elige el umbral por modelo en validación como
% el de mayor recall con precisión >= 0,70. Eso es correcto para comparar
% arquitecturas entre sí (RE2.3), pero el IOV de RE3.1 está comprometido con
% confianza 0,5, NMS IoU 0,3 y verdadero positivo a IoU 0,5. Este capítulo
% reporta en el punto de operación comprometido; si se añade el umbral elegido
% en validación, va como lectura secundaria y rotulada como tal, nunca en su
% lugar.
%
% El protocolo de revisión experta es parte del medio de verificación de RE3.1:
% describir qué se le entrega al revisor (muestra, formato, orden), qué decide y
% cómo se registra su decisión.
% ===================================================================

% ===================================================================
\section{Medición comparativa del tiempo de revisión (RE3.2)}
\label{sec:re32}
% ===================================================================
% TODO Redactar. Es el núcleo del capítulo y la única cifra que responde al
% objetivo general: minutos por hora de audio en condición MANUAL y en condición
% ASISTIDA, medidos sobre material comparable.
%
% El IOV comprometido exige:
%   - minutos por hora de audio en cada condición;
%   - la descripción del material comparable y del orden de revisión (para que el
%     aprendizaje del revisor no se confunda con el efecto de la asistencia);
%   - el tiempo de procesamiento del modelo por hora de audio, que es parte del
%     costo de la condición asistida;
%   - no fijar un porcentaje de ahorro antes de medirlo.
%
% TODO TABLA: condición (manual / asistida) x minutos por hora de audio, con el
%             material, el revisor y el orden en que revisó cada bloque.
% TODO TABLA: costo de inferencia por hora de audio en CPU y en GPU
%             (docs/system_requirements.md ya tiene las mediciones por modelo;
%             aquí van las del modelo seleccionado).
% TODO FIGURA: cobertura (recall por evento a IoU 0,5) frente a propuestas por
%              hora de audio, con el punto de operación comprometido marcado.
%              Es la magnitud que explica POR QUÉ baja el tiempo: cuántas
%              propuestas hay que inspeccionar y cuánto audio queda sin ninguna.
%
% El precedente para el indicador derivado de la salida del sistema (reducción
% de 22 veces del volumen que un operador debe procesar) es
% \citep{Campos_Krofel_Rio-Maior_Renna_2025}; el precedente para el ahorro de
% tiempo humano en primates es \citet{heinicke_2015}. El primero no sustituye a
% la medición cronometrada que este RE compromete: la acompaña.
% ===================================================================

% ===================================================================
\section{Análisis cualitativo de errores y hallazgos (RE3.3)}
\label{sec:re33}
% ===================================================================
% TODO Redactar. Clasificar una muestra de predicciones de ALTA CONFIANZA en las
% tres categorías comprometidas por el IOV: evento no anotado, ruido y error de
% encuadre. La taxonomía geométrica A/B/C/D de la Tabla~\ref{tab:taxonomia-fp} es
% el instrumento con que se separan; el reporte debe poder leerse en las tres
% categorías comprometidas:
%   error de encuadre  <- A (y B, división/fusión, si se declara la equivalencia)
%   evento no anotado  <- D resuelto por el revisor como evento válido
%   ruido              <- D resuelto por el revisor como falso positivo genuino
% Decidir y declarar explícitamente esa correspondencia; C (clase incorrecta) no
% entra en ninguna de las tres y necesita su propio párrafo.
%
% TODO TABLA: conteo y porcentaje por categoría sobre la muestra revisada.
% TODO FIGURA: ejemplos visuales sobre el espectrograma de (a) aciertos,
%              (b) errores comunes y (c) detecciones correctas ausentes de la
%              anotación de referencia. Las tres las pide el IOV.
% TODO FIGURA: distribución de banda de frecuencia, duración y puntuación de las
%              predicciones sin anotación, superpuesta a la de los eventos
%              anotados de cada clase.
%
% CUIDADO CON LA REDACCIÓN: mientras el revisor no resuelva una predicción de
% tipo D, su parecido geométrico con los eventos anotados es INDICIO de una
% omisión en la referencia, no prueba. Con revisión experta -que este RE sí
% compromete- la afirmación pasa a ser una medición; sin ella, se declara como
% limitación.
%
% MATERIAL YA EXISTENTE: research/reportes/RE_3-3_taxonomia-discrepancias.tex y
% su cuaderno reparten las discrepancias A/B/C/D sobre todo el corpus (k-fold en
% train) y dejan la tabla de hallazgos; ahí están las figuras y las cifras que
% esta sección tiene que absorber.
% ===================================================================

% ===================================================================
\section{Discusión}
\label{sec:discusion-oe3}
% ===================================================================
% TODO Redactar siguiendo las pautas, sin introducir resultados nuevos:
%   (1) resumen de los resultados principales, pasando de las cifras a su
%       lectura descriptiva;
%   (2) qué significan: si el tiempo de revisión baja lo suficiente para que el
%       equipo de investigación adopte el flujo asistido, y a costa de qué;
%   (3) consistencia con la literatura del Capítulo~\ref{cap:estado}
%       (Heinicke et al. 2015 y Campos et al. 2025 como puntos de comparación);
%   (4) generalización: qué se sostiene con otro corpus, otro revisor u otras
%       especies, y qué no;
%   (5) limitaciones: tamaño de la muestra revisada, un solo revisor, efecto de
%       aprendizaje y dependencia del punto de operación.
% ===================================================================
